\section*{Contributing Guidelines}
\addcontentsline{toc}{section}{Contributing Guidelines}

\subsection*{Research-Driven Contribution Philosophy}

Symphony's contribution guidelines reflect our commitment to advancing the field of AI-first development environments through rigorous research and collaborative innovation. Contributors participate in both theoretical advancement and practical implementation, ensuring that every contribution enhances our understanding of AI-assisted software development.

Our contribution philosophy emphasizes:
\begin{compactlist}
    \item Empirical validation of architectural decisions and design choices
    \item Theoretical grounding in software engineering and AI research literature
    \item Collaborative peer review and academic-quality documentation
    \item Open sharing of research insights and experimental results
\end{compactlist}

\subsection*{Code Quality and Research Standards}

Code contributions to Symphony must meet both software engineering best practices and academic research standards. This dual requirement ensures that our codebase serves as both a functional system and a research artifact suitable for academic publication and peer review.

\begin{table}[ht]
\centering
\begin{tabular}{@{}lll@{}}
\toprule
\textbf{Quality Dimension} & \textbf{Standard} & \textbf{Validation Method} \\
\midrule
Code Quality & Industry best practices & Automated linting and review \\
Performance & Research benchmarks & Empirical measurement \\
Documentation & Academic standards & Peer review process \\
Testing & Research reproducibility & Complete test coverage \\
Security & Production deployment & Security audit and analysis \\
\bottomrule
\end{tabular}
\caption{Code Quality Standards and Validation Methods}
\end{table}

\subsection*{Commit Conventions and Research Traceability}

Our commit conventions support research traceability and enable systematic analysis of development patterns and architectural evolution. Commit messages follow academic citation standards and provide sufficient context for research analysis.

\begin{alertbox}
Commit messages must include research context and architectural rationale to support future analysis and academic publication of development insights.
\end{alertbox}

Commit message structure:
\begin{compactlist}
    \item \textbf{Type}: Research domain or system component affected
    \item \textbf{Scope}: Specific architectural element or research area
    \item \textbf{Description}: Clear explanation of changes and research implications
    \item \textbf{Rationale}: Theoretical or empirical justification for changes
\end{compactlist}

\subsection*{Pull Request Process and Peer Review}

The pull request process integrates software engineering practices with academic peer review standards. Each pull request undergoes evaluation for both technical correctness and research contribution quality.

Pull request evaluation criteria:
\begin{expandedlist}
    \item \textbf{Technical Correctness}: Code functionality, performance, and security
    \item \textbf{Research Contribution}: Advancement of theoretical understanding or empirical knowledge
    \item \textbf{Documentation Quality}: Academic-standard documentation and explanation
    \item \textbf{Experimental Validation}: Empirical evidence supporting design decisions
    \item \textbf{Integration Impact}: Effects on overall system architecture and research objectives
\end{expandedlist}

\subsection*{Research Documentation Requirements}

All contributions must include documentation that meets academic publication standards. Documentation serves both practical and research purposes, enabling system usage and supporting academic dissemination of research findings.

\begin{infobox}[title=Academic Documentation Standards]
Documentation must be suitable for inclusion in academic publications and conference presentations. This requirement ensures that our research contributions can be effectively shared with the broader research community.
\end{infobox}

Documentation requirements include:
\begin{compactlist}
    \item Theoretical background and related work analysis
    \item Design rationale with empirical or theoretical justification
    \item Implementation details with performance characteristics
    \item Evaluation methodology and experimental results
    \item Implications for future research and development
\end{compactlist}

\subsection*{Community Collaboration and Research Ethics}

Symphony's development community operates according to academic research ethics and collaborative principles. Contributors participate in a research community dedicated to advancing knowledge while maintaining the highest standards of intellectual integrity.

Community guidelines emphasize:
\begin{expandedlist}
    \item \textbf{Open Research}: Transparent sharing of methods, data, and results
    \item \textbf{Collaborative Innovation}: Joint problem-solving and knowledge creation
    \item \textbf{Ethical Conduct}: Responsible research practices and community interaction
    \item \textbf{Knowledge Sharing}: Active contribution to collective understanding
\end{expandedlist}

\subsection*{Recognition and Academic Credit}

Contributors to Symphony receive appropriate academic recognition for their research contributions. Our recognition system acknowledges both technical implementation and theoretical advancement, supporting contributors' academic and professional development.

Recognition mechanisms include:
\begin{compactlist}
    \item Co-authorship on academic publications arising from contributions
    \item Attribution in research presentations and conference talks
    \item Recognition in project documentation and academic acknowledgments
    \item Opportunities for collaborative research and publication
\end{compactlist}
